% ============================================================
\section{Background and Related Work}
\label{sec:background}
% ============================================================

\textbf{Related work: three lines, and the gap between them.}
\emph{(i)~Forgetting} is studied in continual learning, NLP and fine-tuning
dynamics~\cite{kirkpatrick2017ewc, li2017lwf, luo2023forgetting,
kumar2022finetuning, andreassen2023evolution}, and supplies the prescriptions
we test; we use ``forgetting'' throughout in the narrower task-level sense of
degradation relative to zero-shot, not the continual-learning sense of loss on
a prior task.  \emph{(ii)~TSFM adaptation} recipes span probabilistic
fine-tuning, generative pre-training, LLM reprogramming and cross-task
transfer~\cite{rasul2024lagllama, liu2024timer, jin2024timellm,
zhou2023onefitsall}, and are evaluated on target-task accuracy.
\emph{(iii)~Representation analysis} supplies the measurement---representation
similarity and probing~\cite{kornblith2019cka, hewitt2019structural,
pimentel2020information}---including work probing frozen TSFM internals
directly~\cite{wilinski2025representations}, where we instead track how
representations \emph{change under fine-tuning}.  \textbf{None of the three
establishes that the pre-trained representation was worth preserving on the
target task}, which is the premise (i)'s prescriptions rest on and what this
paper measures.

\textbf{Pre-training value gate.}
The gate is deliberately cheap and \emph{downstream-only}: the practitioner
runs the frozen model zero-shot on the target set and compares against a
lookback-96 linear regression \emph{fit on that same dataset}---one ridge-OLS
map $\R^{96 \times F} \to \R^{h \times F}$ estimated on the training split and
applied unchanged to the evaluation windows. Neither pre-training data nor
pre-training compute is required.
We report $R^2_{\text{task}}(\text{PT}) = 1 - \text{MSE}_{\text{ZS}} /
\text{MSE}_{\text{Linear}}$ as a continuous quantity for every cell
(Table~\ref{tab:r2task}) and mark $0.20$ only as an operating point, so
readers can apply alternative thresholds; \S\ref{sec:analysis:prospective}
shows the choice makes no difference to the outcome here.
The threshold is calibrated to ETT short-horizon benchmarks, where linear
baselines are famously competitive with deep
forecasters~\cite{zeng2023transformers}.  That competitiveness is the whole
reason the gate has teeth: \textbf{5 of our 32 screened cells clear it}, all
of them ETTh2 on Moirai, and the remaining 27 are cells in which the
pre-trained model has no measured advantage for fine-tuning to destroy.
Substituting a per-window trend extrapolation---a linear fit to
each evaluation window's own last 96 steps, using no training data---inflates
the denominator enough that a constant predictor beats it on ETTh1 and
Weather, and moves 17 cells across the threshold. An earlier version of this
work used that estimator; Appendix~\ref{app:gatecorrection} documents the
correction and its consequences in full.

\textbf{Task-level consequence: forgetting.}
We define forgetting as signed relative change in zero-shot-canonical
task metric on held-out test:
$\text{forg.\%} = 100 \cdot (M_{\text{FT}} - M_{\text{ZS}}) /
|M_{\text{ZS}}|$ (MSE for forecasting, so positive = worse).
Forg.$>0$ indicates \emph{task-level forgetting}: degraded performance
against the zero-shot baseline.

\textbf{Decoupling harmful forgetting from benign specialisation.}
A reviewer-critical distinction: representational change that redirects
features toward the target (\emph{benign specialisation}) must be
separated from change that discards generally useful structure
(\emph{harmful forgetting}).  An earlier version of this work operationalised the distinction with a
\emph{probe-sign asymmetry} criterion: trained-objective decodability up,
task-orthogonal decodability flat or down.  That criterion is retracted
(Appendix~\ref{sec:retraction}).
What remains is \emph{value-gate conditioning}.  Representational change can be
harmful only where the structure it discards was valuable, i.e.\ in a gate-pass
cell; in a gate-fail cell the same change is benign by construction, because no
demonstrated advantage existed to lose.  Gate-passing plus forgetting${>}0$ is
therefore necessary for harm but not sufficient: the question of whether the
encoder's \emph{change} carried the loss is answered by intervention---comparing
full fine-tuning against a frozen encoder on the same data for the same
epochs---and not by observing the representation.  \S\ref{sec:method} states the
resulting three-clause definition, and \S\ref{sec:analysis:degradation} reports
that no cell we screened satisfies it.  We stress that this is a statement about
what these benchmarks can evidence, not a proof that harmful forgetting never
occurs; \S\ref{sec:limitations} sets out what would be needed to exhibit it.
